\section{Evidence Families}

Across states, certification evidence tends to fall into recurring families.
The exact deadline, official, and audit type varies, but the public proof
problem is stable enough to model.

\begin{center}
\small
\begin{tabularx}{\textwidth}{lX}
\toprule
Evidence family & Public question \\
\midrule
Canvass arithmetic & Do reporting-unit, contest, batch, and jurisdiction totals agree after official corrections? \\
Lineage & Are precinct splits, merges, renames, and boundary changes explicit across cycles? \\
Source integrity & Do public exports, summaries, CVRs, and manifests still hash to the declared package evidence? \\
Ballot accounting & Do batches, manifests, accepted ballots, rejected ballots, and counted ballots reconcile? \\
CVR reconciliation & Do ballot-level CVR rows reproduce published summary totals where CVRs are available? \\
Audit completion & Did the required hand audit, RLA, or recount produce an auditable pass/escalate record? \\
Voter-facing inclusion & Can a voter check inclusion without creating a receipt for candidate choices? \\
District aggregation & Do district totals follow from a count package and a plan/context hash? \\
\bottomrule
\end{tabularx}
\end{center}

Risk-limiting audits deserve special care because they are statistical and
jurisdiction-specific. NCSL describes RLAs as audits that limit the risk of
certifying an incorrect outcome and notes that audit completion may be tied to
canvass or certification timelines \citep{ncsl-post-election-audits}. RCOUNT
therefore records replayable samples, observations, margins, stopping methods,
and adapter metadata, but leaves statutory adequacy to jurisdiction-specific
rules.
